\documentclass[../main.tex]{subfiles}
\begin{document}

\section{Lecture 23 -- Fluid Mechanics from the Action Principle}\label{sec:lec23}
\begin{center}
  \textbf{Date:} June 24, 2026
\end{center}

\subsection*{Overview}
The final lecture returns to a more structural question.  Throughout the course
we mostly wrote the local equations of motion directly: force balance in
elasticity, continuity plus Euler's equation for ideal fluids, and
Navier--Stokes for viscous fluids.  The viewpoint of this lecture is different:
ask whether the equations can be obtained from one scalar object, the
\textbf{action}.

In ordinary analytical mechanics a trajectory $q_i(t)$ is physical when it makes

action stationary.  The lecture generalizes this statement to fields
$\phi_a(\bvec x,t)$, whose degrees of freedom are labelled by both space and
time.  The payoff is conceptual: if a theory has an action, then continuous
symmetries of the action imply conservation laws by Noether's theorem.
For ideal fluids this gives a unified origin for energy conservation, momentum
conservation, and Kelvin circulation conservation.  Viscosity is different:
like friction in mechanics, it is dissipative and is not naturally obtained from
a standard action principle.

\medskip\noindent\textbf{Topics:}
\begin{enumerate}
  \item Reminder: action, Euler--Lagrange equations, and Noether's theorem in
    mechanics.
  \item Local field actions and the field Euler--Lagrange equations.
  \item Full variational derivation of the field Euler--Lagrange equations.
  \item Example: vibrating string and the wave equation.
  \item Functional derivative viewpoint: why
    $\delta\phi(x)/\delta\phi(y)=\delta(x-y)$.
  \item Noether theorem for fields and conserved currents.
  \item Material/Lagrangian-coordinate action for ideal fluids.
  \item Symmetries behind the conservation laws of ideal fluids.
  \item Incompressible ideal fluid action and pressure as a Lagrange multiplier.
\end{enumerate}

\subsection{From Mechanics to Fields}

\subsubsection*{Ordinary analytical mechanics}
For finitely many generalized coordinates $q_i(t)$, the action is
\begin{equation}
  S[q]=\int_{t_1}^{t_2} L(q_i,\dot q_i,t)\,dt .
  \label{eq:lec23_mechanical_action}
\end{equation}
Stationarity, $\delta S=0$ for arbitrary variations $\delta q_i(t)$ that vanish
at the endpoints, gives
\begin{equation}
  \boxed{\pdv{L}{q_i}-\dv{t}\pdv{L}{\dot q_i}=0 .}
  \label{eq:lec23_mechanical_el}
\end{equation}
Thus one scalar function $L$ produces all equations of motion.

Noether's theorem says more.  If a continuous transformation with infinitesimal
parameter $s$ changes the Lagrangian only by a total derivative,
\begin{equation}
  \delta L=\dv{\Lambda}{t},
  \label{eq:lec23_mech_noether_condition}
\end{equation}
then the corresponding quantity
\begin{equation}
  Q=\pdv{L}{\dot q_i}\,\delta q_i-\Lambda
  \label{eq:lec23_mech_noether_charge}
\end{equation}
is conserved on solutions.  The subtraction of $\Lambda$ is important.  The
standard energy conservation law from time translations is the famous case in
which this total-derivative term matters.

\subsubsection*{Fields as infinitely many degrees of freedom}
A field theory replaces the finite list $q_i(t)$ by fields
\begin{equation}
  \phi_a=\phi_a(x^\mu),\qquad x^\mu=(t,x,y,z),
\end{equation}
where $a$ labels the different fields.  A mechanical coordinate is a special
case of a field depending only on time.  The local field-theory action is
\begin{equation}
  \boxed{S[\phi]=\int d^4x\,\mathcal L\!\left(\phi_a,\partial_\mu\phi_a\right)}
  \label{eq:lec23_field_action}
\end{equation}
with Lagrangian density $\mathcal L$.  ``Local'' means that at a point $x$ the
density depends only on the fields and their derivatives at the same point, not
on values at far-away points.

\paragraph{Physical intuition for locality.}
A field removes mysterious action at a distance.  For example, instead of saying
that Earth pulls the Moon directly across empty space, one may say: Earth
creates a gravitational potential field $\Phi$, the field obeys a local
differential equation, and the Moon responds to the local gradient
$-\grad\Phi$ at its position.  Information is mediated by fields and local
equations.

\paragraph{Dimensions.}
In mechanics, $[L]=\text{energy}$ and $[S]=\text{action}$.  In field theory,
$S$ still has dimensions of action, but because $S=\int dt\,d^3x\,\mathcal L$,
\begin{equation}
  [\mathcal L]=\frac{\text{energy}}{\text{volume}}.
\end{equation}
So $\mathcal L$ is an energy density rather than an energy.

\subsection{Derivation of the Field Euler--Lagrange Equations}
Consider one field $\phi$ first.  The multi-field case just adds an index $a$.
Vary the field by an arbitrary small function that vanishes on the spacetime
boundary,
\begin{equation}
  \phi(x)\longrightarrow \phi(x)+\delta\phi(x),
  \qquad \delta\phi\big|_{\partial\Omega}=0 .
\end{equation}
The first-order change in the action is
\begin{align}
  \delta S
  &=\int_\Omega d^4x\,\delta\mathcal L  \\
  &=\int_\Omega d^4x\left[
    \pdv{\mathcal L}{\phi}\,\delta\phi
    +\pdv{\mathcal L}{(\partial_\mu\phi)}\,\delta(\partial_\mu\phi)
  \right].
  \label{eq:lec23_delta_s_chain}
\end{align}
Variation and differentiation commute, so
\begin{equation}
  \delta(\partial_\mu\phi)=\partial_\mu(\delta\phi).
\end{equation}
Now integrate the second term by parts:
\begin{align}
  \int_\Omega d^4x\,
  \pdv{\mathcal L}{(\partial_\mu\phi)}\partial_\mu(\delta\phi)
  &=\int_\Omega d^4x\,\partial_\mu\left(
    \pdv{\mathcal L}{(\partial_\mu\phi)}\delta\phi\right) \\
  &\quad-\int_\Omega d^4x\,
    \partial_\mu\left[\pdv{\mathcal L}{(\partial_\mu\phi)}\right]\delta\phi .
\end{align}
By Gauss' theorem the first term is a boundary flux.  It vanishes because
$\delta\phi=0$ on the boundary.  Therefore
\begin{equation}
  \delta S=\int_\Omega d^4x\left[
    \pdv{\mathcal L}{\phi}
    -\partial_\mu\left(\pdv{\mathcal L}{(\partial_\mu\phi)}\right)
  \right]\delta\phi .
  \label{eq:lec23_delta_s_final}
\end{equation}
For this to vanish for every possible interior variation $\delta\phi(x)$, the
coefficient must vanish pointwise.  Hence
\begin{equation}
  \boxed{
  \pdv{\mathcal L}{\phi_a}
  -\partial_\mu\left(\pdv{\mathcal L}{(\partial_\mu\phi_a)}\right)=0
  }
  \label{eq:lec23_field_el}
\end{equation}
for each field $\phi_a$.

\nt{The field Euler--Lagrange equation is local: it is a differential equation
involving the field and its derivatives at the same point.  This locality came
straight from the local form of the action.}

\subsection{Example: Vibrating String}

\subsubsection*{Physical setup}
Let $\phi(x,t)$ be the transverse displacement of a stretched string.  Let
$\rho$ be the mass per unit length and $T$ the tension.  For small slopes,
\begin{equation}
  \mathcal L=\frac12\rho(\partial_t\phi)^2-\frac12T(\partial_x\phi)^2 .
  \label{eq:lec23_string_lagrangian}
\end{equation}
This is exactly ``kinetic energy density minus potential energy density.''  The
potential energy comes from stretching: if neighboring points have different
transverse displacements, $\partial_x\phi\ne0$, the string length increases and
tension stores energy.

\subsubsection*{Derivation}
The density does not depend explicitly on $\phi$, only on its derivatives, so
\begin{equation}
  \pdv{\mathcal L}{\phi}=0,
  \qquad
  \pdv{\mathcal L}{(\partial_t\phi)}=\rho\partial_t\phi,
  \qquad
  \pdv{\mathcal L}{(\partial_x\phi)}=-T\partial_x\phi .
\end{equation}
Substituting into \eqref{eq:lec23_field_el},
\begin{align}
  0-\partial_t(\rho\partial_t\phi)-\partial_x(-T\partial_x\phi)&=0,\\
  \rho\partial_t^2\phi-T\partial_x^2\phi&=0.
\end{align}
Thus
\begin{equation}
  \boxed{\partial_t^2\phi=c^2\partial_x^2\phi,
  \qquad c^2=\frac{T}{\rho}.}
  \label{eq:lec23_string_wave}
\end{equation}
The wave speed increases with tension and decreases with inertia per unit
length, exactly as physical intuition suggests.

\subsection{Functional Derivative Viewpoint}
The field $\phi(x)$ may be thought of as infinitely many variables: one variable
for each spatial point.  In a discretized version, write
\begin{equation}
  \phi_i\equiv \phi(x_i).
\end{equation}
Ordinary independent variables satisfy
\begin{equation}
  \pdv{\phi_i}{\phi_j}=\delta_{ij}.
\end{equation}
In the continuum limit the Kronecker delta becomes the Dirac delta:
\begin{equation}
  \boxed{\frac{\delta\phi(y)}{\delta\phi(x)}=\delta^{(d)}(y-x).}
  \label{eq:lec23_functional_delta}
\end{equation}
The intuition is direct: changing the field at $x$ does not change the field at
$y$ unless $x=y$.

Using this rule, compute the functional derivative of
$S=\int d^dy\,\mathcal L(\phi(y),\partial_\mu\phi(y))$:
\begin{align}
  \frac{\delta S}{\delta\phi(x)}
  &=\int d^dy\left[
    \pdv{\mathcal L}{\phi(y)}\frac{\delta\phi(y)}{\delta\phi(x)}
    +\pdv{\mathcal L}{(\partial_\mu\phi(y))}
      \frac{\delta(\partial_\mu\phi(y))}{\delta\phi(x)}
  \right] \\
  &=\int d^dy\left[
    \pdv{\mathcal L}{\phi(y)}\delta(y-x)
    +\pdv{\mathcal L}{(\partial_\mu\phi(y))}\partial_{y^\mu}\delta(y-x)
  \right].
\end{align}
The second term is integrated by parts in $y$:
\begin{equation}
  \int d^dy\,F(y)\partial_{y^\mu}\delta(y-x)
  =-\partial_{x^\mu}F(x).
\end{equation}
Therefore
\begin{equation}
  \boxed{
  \frac{\delta S}{\delta\phi(x)}
  =\pdv{\mathcal L}{\phi}(x)
  -\partial_\mu\left[\pdv{\mathcal L}{(\partial_\mu\phi)}\right](x).}
  \label{eq:lec23_functional_derivative_el}
\end{equation}
Stationarity is simply
\begin{equation}
  \frac{\delta S}{\delta\phi_a(x)}=0
  \qquad\text{for every field and every point.}
\end{equation}

\subsection{Noether Theorem for Fields}
Suppose a continuous infinitesimal transformation changes the fields by
\begin{equation}
  \phi_a\longrightarrow \phi_a+s\,\delta\phi_a+O(s^2)
\end{equation}
and changes the Lagrangian density by a total divergence,
\begin{equation}
  \delta\mathcal L=\partial_\mu\Lambda^\mu .
  \label{eq:lec23_field_noether_condition}
\end{equation}
Then the Noether current is
\begin{equation}
  \boxed{j^\mu=\frac{\partial\mathcal L}{\partial(\partial_\mu\phi_a)}\,
  \delta\phi_a-\Lambda^\mu,}
  \label{eq:lec23_noether_current}
\end{equation}
with summation over $a$, and on solutions it obeys
\begin{equation}
  \boxed{\partial_\mu j^\mu=0.}
  \label{eq:lec23_current_conservation}
\end{equation}

In ordinary mechanics the conserved object is a number.  In field theory the
conserved object is a \textbf{current}.  Writing $j^\mu=(\rho_Q,\bvec j_Q)$,
\begin{equation}
  \partial_t\rho_Q+\divg\bvec j_Q=0.
\end{equation}
The amount of $Q$ inside a volume changes only because $Q$ flows through the
boundary.  This is the same structure as charge conservation,
$\partial_t\rho_e+\divg\bvec j_e=0$.

\subsection{Ideal Fluids from an Action}

\subsubsection*{Why use material coordinates?}
The Eulerian velocity field $\vel(\bvec x,t)$ is the natural language for many
fluid calculations, but the action principle is simplest in the material
(Lagrangian) description.  Label each fluid element by
\begin{equation}
  \refcoord=(X^1,X^2,X^3),
\end{equation}
often chosen to be its initial position.  The dynamical fields are the current
positions of the labelled elements,
\begin{equation}
  x^i=x^i(\refcoord,t),\qquad i=1,2,3.
\end{equation}
The velocity of that element is $\dot x^i(\refcoord,t)$.

Let
\begin{equation}
  J=\det\left(\pdv{x^i}{X^A}\right)
  \label{eq:lec23_jacobian}
\end{equation}
be the Jacobian from label space to physical space.  Since mass is conserved,
\begin{equation}
  \rho(\bvec x,t)=\frac{\rho_0(\refcoord)}{J},
  \label{eq:lec23_mass_jacobian}
\end{equation}
where $\rho_0(\refcoord)$ is the mass density in label space.

\subsubsection*{Compressible ideal fluid action}
For a compressible ideal fluid, a standard material-coordinate action is
\begin{equation}
  \boxed{
  S=\int dt\,d^3X\,\rho_0(\refcoord)
  \left[\frac12\dot{\bvec x}^{\,2}-\epsilon(\rho)-\Phi(\bvec x)\right].}
  \label{eq:lec23_compressible_fluid_action}
\end{equation}
Here $\epsilon$ is the internal energy per unit mass and $\Phi$ is an external
potential per unit mass, for example $\Phi=gz$ in uniform gravity.  The pressure
is hidden in the dependence of $\epsilon$ on density, and density is hidden in
$J$ via \eqref{eq:lec23_mass_jacobian}.  Varying $x^i(\refcoord,t)$ gives the
Euler equation of ideal-fluid dynamics.  The algebra is longer than in the
string example because varying $x^i$ also varies the determinant $J$.

The useful geometric identity mentioned in the lecture is the Piola identity,
\begin{equation}
  \pdv{}{X^A}\left(J\pdv{X^A}{x^i}\right)=0,
  \label{eq:lec23_piola_identity}
\end{equation}
which is one of the tools that turns the variation of the determinant into the
pressure-gradient force in Euler's equation.

\subsubsection*{Physical meaning}
The action is still kinetic minus potential energy.  The kinetic part tracks the
motion of each labelled fluid element.  The internal-energy part penalizes
compression or expansion because changing $J$ changes $\rho$.  The external
potential part gives body forces such as gravity.  Therefore the usual pressure
force in Euler's equation is not put in by hand; it arises because neighboring
fluid elements resist changes of local volume.

\subsection{Conservation Laws of an Ideal Fluid as Symmetries}
Ideal fluids have three central conservation laws, and the action viewpoint
organizes them by their symmetries.

\begin{center}
\begin{tabular}{lll}
\toprule
Conservation law & Symmetry & Comment \\
\midrule
Energy & time translations & gives energy density and energy flux \\
Momentum & spatial translations & also survives for viscous flow without external forces \\
Kelvin circulation/vorticity & relabelling of fluid elements & special to ideal, non-dissipative flow \\
\bottomrule
\end{tabular}
\end{center}

\paragraph{Energy.}
If the action has no explicit time dependence, time translation symmetry implies
energy conservation.  In field language this is not just one number but a local
continuity equation for energy density and energy flux.

\paragraph{Momentum.}
If there is no preferred point in space, spatial translation symmetry implies
momentum conservation.  Even when viscosity is present and a standard action is
lost, total momentum is still conserved if all forces are internal and no
external body force acts.

\paragraph{Kelvin circulation.}
The labels $\refcoord$ are names of fluid elements.  Physics cannot depend on
which names we choose.  This relabelling freedom is a symmetry of the ideal
fluid action.  The corresponding Noether statement is Kelvin's theorem: the
circulation around a material loop is conserved,
\begin{equation}
  \dv{}{t}\oint_{C(t)}\vel\cdot d\bvec l=0,
\end{equation}
under the usual ideal/barotropic assumptions.

\subsection{Incompressible Ideal Fluid and Pressure as a Constraint}
For an incompressible fluid the density is fixed.  A material volume element
must keep its volume, so
\begin{equation}
  J=1.
  \label{eq:lec23_incompressibility_j}
\end{equation}
The action can impose this with a Lagrange multiplier $p(\refcoord,t)$:
\begin{equation}
  \boxed{
  S_{\rm inc}=\int dt\,d^3X\left\{
  \rho_0\left[\frac12\dot{\bvec x}^{\,2}-\Phi(\bvec x)\right]
  +p(\refcoord,t)\,[J-1]
  \right\}.}
  \label{eq:lec23_incompressible_action}
\end{equation}
Varying with respect to $p$ gives $J=1$.  Varying with respect to $\bvec x$ gives
the incompressible Euler equation, with $p$ becoming the physical pressure.  This
explains a central fact about incompressible flow: pressure is not supplied by
an equation of state.  It is the constraint force that enforces
$\divg\vel=0$.

\subsection*{Key Takeaways}
\begin{itemize}
  \item A field theory action is an integral of a local Lagrangian density over
    spacetime.
  \item Stationarity gives local Euler--Lagrange differential equations.
  \item Functional derivatives are ordinary derivatives with respect to an
    infinite set of variables; the Dirac delta replaces the Kronecker delta.
  \item Noether's theorem in field theory gives conserved currents
    $\partial_\mu j^\mu=0$, not merely conserved numbers.
  \item Ideal fluids admit an action in material coordinates.  Viscous fluids do
    not, because viscosity is dissipative like friction.
  \item Incompressibility is the constraint $J=1$, and pressure is the Lagrange
    multiplier enforcing it.
\end{itemize}

\end{document}
